\documentclass[11pt,a4paper]{article}

\usepackage[utf8]{inputenc}
\usepackage[T1]{fontenc}
\usepackage[english]{babel}
\usepackage[margin=2.2cm]{geometry}
\usepackage{array}
\usepackage{booktabs}
\usepackage{longtable}
\usepackage{enumitem}
\usepackage{xcolor}
\usepackage{microtype}
\usepackage[hidelinks]{hyperref}

% Items that need a decision or information, printed in red.
\newcommand{\ph}[1]{\textcolor{red!65!black}{[#1]}}
\newcolumntype{L}[1]{>{\raggedright\arraybackslash}p{#1}}

\setlength{\parindent}{0pt}
\setlength{\parskip}{5pt}
\setlist{itemsep=3pt, topsep=3pt}
\renewcommand{\arraystretch}{1.25}

\title{Framework paper: proposed outline for approval}
\author{NISC Lab, Cognitive Science Department, University of Messina}
\date{14 September 2026}

\begin{document}
\maketitle

\section*{What is asked}

Approval of the scope, framing and structure below, before the paper
is rewritten into them. The structure follows the earlier laboratory
manuscript (Daghooghi et al.), and every point its reviewers raised has
a place in it (Section~\ref{sec:reviewers}). Items in \ph{red brackets}
need a decision or information from the person named. Technical team:
P.~Kazemi, E.~Zandi, B.~Baharizadeh. Cognitive team: P.~Daghooghi,
Y.~Asghari Dizaj, H.~Khorshidi, E.~Bakhtiari.

% ,,,,,,,,,,,,,,,,,,,,,,
\section{Decisions and information needed}
\label{sec:decisions}
% ,,,,,,,,,,,,,,,,,,,,,,

\begin{enumerate}
    \item \textbf{Ethics approval.} Name of the committee, protocol
    number and date of approval, and confirmation that it covers the
    sessions recorded since July 2026 (33 runs from 24 participants).
    \ph{Prof.~Grasso, E.~Bakhtiari}

    \item \textbf{Title.} Provisional: \emph{Design of an Integrated
    VR-Biofeedback Framework for Controlled Phobia Exposure}.
    \ph{administration}

    \item \textbf{Authors.} Spelling of Haleh (or Halleh) Khorshidi and
    of Febronia Riggio; affiliations of F.~Riggio and C.~M.~Porto.
    \ph{Prof.~Grasso}

    \item \textbf{Venue and template} of the full paper, which set its
    length. \ph{Prof.~Grasso}

    \item \textbf{Origin of the index coefficients.} The weights 0.5,
    0.3 and 0.2 and the three-band rule come from the laboratory's own
    earlier design and appear in no cited source. Proposed: describe
    them as coefficients fixed in advance by the laboratory, with their
    rationale, or cite an internal report if one can be made citable.
    \ph{Prof.~Grasso}

    \item \textbf{Control group.} Proposed: stated as a limitation of a
    methodological study. A later comparison could use a yoked
    condition, in which a participant receives without adaptation the
    stimulus sequence logged for another participant; the command log
    already holds such sequences. \ph{Prof.~Grasso}

    \item \textbf{Agreement between the two extraction
    implementations.} Proposed for a longer version only. On the two
    current captures they agree on RMSSD ($r = 0.72$ and $0.88$) but
    not on heart rate ($r = -0.12$ and $-0.05$), because the
    manufacturer's detector marks 14 to 24\% extra beats; this needs
    checking before it is published. \ph{Prof.~Grasso}

    \item \textbf{Environment parameters:} starting height, height
    change per command, lower and upper limits, and run duration in
    each session. \ph{environment developers}

    \item \textbf{Participants and questionnaires:} recruitment,
    inclusion cut-off and the scale it refers to, exclusion criteria,
    instruments before and after exposure, target sample.
    \ph{cognitive team}
\end{enumerate}

% ,,,,,,,,,,,,,,,,,,,,,,
\section{Scope and framing}
\label{sec:scope}
% ,,,,,,,,,,,,,,,,,,,,,,

\begin{itemize}
    \item \textbf{A framework paper.} The contribution is the apparatus:
    sensing, calibration, index, control law, safeguards and session
    record, specified closely enough to be reproduced. Data collection
    is in progress and the paper claims no outcomes, which answers the
    objection that the earlier manuscript was submitted before its study
    was complete.

    \item \textbf{Results limited to operation:} a summary of the
    recorded sessions and one illustrative session. Questionnaire
    outcomes are reported later, in separate papers.

    \item \textbf{Experimental wording} throughout (exposure,
    participant, operator, target band), as the laboratory decided.

    \item \textbf{A scenario-independent description,} with the height
    environment as the worked example because it is the environment
    with recorded sessions. Further environments are described once
    they have recorded sessions of their own.

    \item \textbf{The current system only.} The earlier manuscript
    described the previous version; Section~\ref{sec:changes} lists
    what has changed since.

    \item \textbf{Length.} The present methodology, about 9,700 words,
    is kept as the long version. The conference text is condensed from
    it, with each fact stated once.
\end{itemize}

% ,,,,,,,,,,,,,,,,,,,,,,
\section{Proposed structure}
\label{sec:structure}
% ,,,,,,,,,,,,,,,,,,,,,,

Section numbers under \emph{Source} refer to the current methodology
draft.

\begin{longtable}{@{}L{3.4cm}L{8.6cm}L{3.6cm}@{}}
\toprule
\textbf{Section} & \textbf{Content} & \textbf{Source and owner} \\
\midrule
\endhead
\bottomrule
\endfoot
Abstract &
    Aim, the framework, how the loop works, state of data collection;
    the phobia term defined at first use. &
    co-authors \\
1 Introduction &
    Specific phobias and exposure. Fear of heights distinguished from
    diagnosed phobia, with prevalence given for each. Limits of in vivo
    exposure, with examples. Why exposure should follow physiology. &
    cognitive team \\
1.1 Related work &
    Virtual exposure with scripted or operator-graded stimuli, with
    examples; systems that adapt to physiology; what a separable, fully
    specified framework adds. &
    cognitive team; shared bibliography (67 entries) \\
\addlinespace[3pt]
\multicolumn{3}{@{}l}{\textbf{2 Methodology}} \\
2.1 Study design &
    Methodological single-group design; session plan; what is evaluated
    and what is not; absence of a control condition stated. &
    cognitive team, with 3.1 \\
2.2 Participants &
    Recruitment; inclusion by a cut-off score on a named scale;
    exclusion criteria; target sample. &
    cognitive team \\
2.3 Materials and apparatus &
    Headset and rendering; sensor unit, chest band, electrode sites and
    skin preparation; software and versions. &
    3.3, condensed \\
\addlinespace[3pt]
\multicolumn{3}{@{}l}{\textbf{3 Adaptive exposure system}} \\
3.1 System overview &
    Closed loop; measurement server and environment client;
    scenario-neutral interface; topology figure. &
    3.1, 3.2, 3.9 \\
3.2 Physiological measures and arousal index &
    Heart rate, RMSSD and phasic conductance with their windows and
    artifact handling; 120~s calibration and its validity check;
    standardised deviations, fixed weights, 3~s smoothing. &
    3.4 to 3.7 \\
3.3 Real-time adaptation &
    Thresholds from each participant's resting distribution; control
    law (increase, hold, decrease); settling gate and rate limit;
    environment parameters in numbers. &
    3.8, 3.10; environment developers \\
3.4 Monitoring, safety and session data &
    Who is present; termination by participant and by operator;
    automatic safeguards and loss of signal; operator console; recorded
    files; data protection. &
    3.11 to 3.13, 3.15 \\
\addlinespace[3pt]
\multicolumn{3}{@{}l}{\textbf{4 Experimental procedure}} \\
4.1 Before exposure &
    Consent, screening and baseline questionnaires. &
    cognitive team \\
4.2 Exposure session &
    Preparation, settling, calibration, run and termination, with
    durations. &
    3.11 \\
4.3 After exposure &
    Questionnaires after the final session. &
    cognitive team \\
\addlinespace[3pt]
5 Results &
    Summary of the recorded sessions, 33 runs from 24 participants so
    far (completed calibrations, run durations, commands sent, episodes
    of signal loss), and one illustrative session as a figure. No
    outcome measures. &
    technical team \\
6 Discussion &
    What the framework makes possible; limitations: single group, index
    not validated against self-report, coefficients fixed in advance. &
    co-authors \\
7 Conclusion &
    The contribution in one paragraph; next steps. &
    co-authors \\
Ethics statement &
    Committee, protocol number, informed consent. &
    \ph{administration} \\
\end{longtable}

Section 3.14 of the draft (validation and reproducibility) and the
comparison of the two extraction implementations stay in the long
version.

% ,,,,,,,,,,,,,,,,,,,,,,
\section{Reviewer points and where they are answered}
\label{sec:reviewers}
% ,,,,,,,,,,,,,,,,,,,,,,

R1 and R2 denote the two reviewers. \emph{Written} means the current
methodology draft already answers the point.

\begin{longtable}{@{}L{5.6cm}L{7.4cm}L{2.6cm}@{}}
\toprule
\textbf{Point raised} & \textbf{Where and how} & \textbf{Status} \\
\midrule
\endhead
\bottomrule
\endfoot
Study still running; outcomes missing; submission premature (R1, R2). &
    Framed as a framework paper with no outcome claims; results limited
    to operation (Sections~\ref{sec:scope} and~\ref{sec:structure}). &
    proposed \\
Prevalence: fear of heights versus diagnosed phobia (R1). &
    Introduction, with a separate figure for each. &
    cognitive team \\
Cognitive aspects and causes of the fear (R1). &
    Introduction, briefly. &
    cognitive team \\
Examples of the limits of in vivo exposure (R1). &
    Introduction. &
    cognitive team \\
Examples of static, non-adaptive virtual exposure systems (R1). &
    Related work, from the shared bibliography. &
    cognitive team \\
Preference for virtual over in vivo exposure (Garc\'{\i}a-Palacios
et al., 2001) and comparative effectiveness (Kuleli et al., 2025) (R1). &
    Introduction; both are already in the shared bibliography. &
    cognitive team \\
Subjective units of distress to be explained (R1). &
    Procedure, if the scale is used. &
    cognitive team \\
Heights used in each session (R1). &
    3.3, in metres: starting height, change per command, limits. &
    \ph{environment developers} \\
Whether the balloon keeps rising while arousal stays low (R1). &
    3.3: control law stated explicitly. While the index stays in the
    lowest band an increase is sent at most once per second, within the
    environment's limits. &
    written \\
Inclusion criteria: cut-off on a validated scale (R1). &
    2.2. &
    cognitive team \\
Ethical approval (R1). &
    Ethics statement with committee and protocol number. &
    \ph{reference needed} \\
How the sensors are attached (R1). &
    2.3: chest band, Lead I electrode sites, finger electrodes on the
    non-dominant hand, skin preparation. &
    written \\
Information repeated (R1). &
    Each fact stated once; condensed text. &
    proposed \\
Who is present during exposure (R1). &
    3.4: an operator at the console and a second researcher for the
    environment and the headset, both in the room throughout. &
    written; roles and training \ph{to confirm} \\
Safety precautions for distress or extreme heart rate (R1). &
    3.4: the participant may stop at any time; the operator stops at any
    sign of distress; implausible heart rate and loss of signal are
    handled automatically; stimulus is withdrawn when arousal is high. &
    written \\
How the target zone is set, and whether it rests on earlier studies
(R1). &
    3.2 and 3.3: set for each participant at the calibration mean plus
    1.28 and 2.33 standard deviations (90th and 99th percentiles). &
    written; coefficients \ph{decision 5} \\
Theoretical framework and literature (R2). &
    Introduction and related work, from the shared bibliography. &
    cognitive team \\
No control group (R2). &
    Stated as a limitation; yoked comparison as an option. &
    \ph{decision 6} \\
Phobia term to be defined in the abstract (R2). &
    Abstract. &
    co-authors \\
Paragraph too long (R2). &
    Shorter paragraphs throughout. &
    all \\
\end{longtable}

% ,,,,,,,,,,,,,,,,,,,,,,
\section{Changes from the system in the earlier manuscript}
\label{sec:changes}
% ,,,,,,,,,,,,,,,,,,,,,,

\begin{itemize}
    \item Conductance enters as its phasic component in microsiemens,
    no longer as a percentage change of the raw level, which drifts
    over minutes.
    \item Each measure is standardised by the participant's resting
    dispersion before weighting, so no measure dominates by scale.
    \item Thresholds sit at the calibration mean plus 1.28 and 2.33
    standard deviations of the index, in place of 1.33 and 2.28 times
    the baseline standard deviation.
    \item The index is smoothed by a 3~s trailing mean.
    \item A calibration that fails a validity check blocks the run.
    \item Commands are held at the start of a run, limited to one per
    second, and withheld while the sensor signal is lost.
    \item Two independent extraction implementations can be selected,
    and archived recordings can be replayed through the live pipeline.
    \item Records carry a participant code; names are no longer stored.
\end{itemize}

% ,,,,,,,,,,,,,,,,,,,,,,
\section{Timeline}
% ,,,,,,,,,,,,,,,,,,,,,,

\begin{itemize}
    \item \textbf{HCI International 2027}, Berlin, 25 to 30 July 2027:
    800-word proposal due 9 October 2026 (draft attached separately),
    notification 20 November 2026, camera-ready paper of 10 to 20 pages
    in the Springer template due 29 January 2027.
    \item \textbf{Full paper:} rewritten into the approved structure one
    section at a time; venue \ph{decision 4}.
\end{itemize}

\end{document}
